\section{Hamiltonian reduction and comparison of complexes}

The constraint $e=1$ for $\mathfrak{sl}_2$ gives a reduction whose maps can be written on every polynomial.
It separates Hamiltonian reduction from the Virasoro charge and from chiral collision operations.
All complexes in the polynomial calculation are algebraic cochain complexes over $\mathbb C$.
Their products use the Koszul sign $(-1)^{|x||y|}$.

\subsection{The polynomial reduction}

Let $e,h,f$ be the linear coordinate functions on $\mathfrak{sl}_2^*$, with Poisson brackets
\[
 \{h,e\}=2e,\qquad \{h,f\}=-2f,\qquad \{e,f\}=h.
\]
Introduce odd generators $b,c$ of cohomological degrees $-1,1$.
The even generators have degree zero.
On the graded commutative algebra
\[
 A=\mathbb C[e,h,f]\otimes\Lambda(b,c)
\]
define a degree-one derivation $D$ by
\begin{equation}\label{eq:cbc26-ds-differential}
 De=0,\qquad Dh=-2ec,\qquad Df=hc,
 \qquad Db=e-1,\qquad Dc=0.
\end{equation}
Thus the Koszul differential imposes $e=1$, and the other term is the Chevalley--Eilenberg differential for the Hamiltonian vector field $\{e,-\}$.
The two terms anticommute.
Direct substitution on the five generators gives $D^2=0$.
Let $R=\mathbb C[u]$ have zero differential and augmentation $u\mapsto0$.
Give $A$ the augmentation
\[
 \varepsilon(e)=1,\qquad
 \varepsilon(h)=\varepsilon(f)=\varepsilon(b)=\varepsilon(c)=0.
\]
These choices concern the displayed differential graded commutative algebra.
They do not give an augmentation of an affine or ghost vertex algebra.

\begin{theorem}[Polynomial reduction and its bar comparison]
\label{thm:cbc26-polynomial-ds}
The maps of augmented differential graded algebras
\begin{align}
 i:R&\longrightarrow A,& i(u)&=ef+\tfrac14h^2,\label{eq:cbc26-ds-inclusion}\\
 p:A&\longrightarrow R,&
 (p(e),p(h),p(f),p(b),p(c))&=(1,0,u,0,0)
 \label{eq:cbc26-ds-projection}
\end{align}
satisfy $pi=1$ and are quasi-isomorphisms.
In particular,
\[
 H^j(A,D)=
 \begin{cases}\mathbb C[u],&j=0,\\0,&j\ne0.\end{cases}
\]
The induced maps $B(i)$ and $B(p)$ on the direct-sum associative reduced bar constructions are quasi-isomorphisms, and $B(p)B(i)=1$.
\end{theorem}

\begin{proof}
The element $K=ef+h^2/4$ satisfies
$DK=ehc-ehc=0$.
The projection kills the image under $D$ of every generator, including $Db=e-1$.
This proves the chain identities and $pi=1$.

Resolve the differential as
\[
 D=(e-1)\partial_b+c(-2e\partial_h+h\partial_f).
\]
Filter by the number of $c$ factors.
The filtration has two steps, so its spectral sequence has no convergence issue.
Its first cohomology calculation is the two-term Koszul complex for $e-1$.
Multiplication by $e-1$ is injective on the polynomial ring.
The resulting complex is
\[
 \mathbb C[h,f]\xrightarrow{\,-2\partial_h+h\partial_f\,}
 c\mathbb C[h,f].
\]
Set $u=f+h^2/4$.
This is a polynomial change of coordinates from $(h,f)$ to $(h,u)$.
In these coordinates the displayed derivative is $-2\partial_h$.
Its kernel is $\mathbb C[u]$, and it is surjective in characteristic zero.
For example, the primitive of $c h^a u^j$ is $-h^{a+1}u^j/(2a+2)$.
Thus cohomology is concentrated in degree zero, and $p$ induces its displayed identification.
The identity $pi=1$ proves the same assertion for $i$.

Here are the bar conventions and the induced maps.
For an augmented differential graded algebra $E$, put $\bar E=\ker\varepsilon$ and
\[
 B(E)=\bigoplus_{n\geq0}(\bar E[1])^{\otimes n}.
\]
Write $sa$ for $a$ considered in $\bar E[1]$, so $|sa|=|a|-1$.
The bar coderivation has components
\[
 b_1(sa)=-s(d_Ea),\qquad
 b_2(sa\otimes sb)=(-1)^{|a|}s(ab).
\]
Extend each component to consecutive positions with the Koszul sign for a degree-one operator.
The Leibniz rule and associativity give $b^2=0$.
For an augmented chain algebra map $q$, set
\[
 B(q)(sa_1\otimes\cdots\otimes sa_n)
 =s q(a_1)\otimes\cdots\otimes s q(a_n).
\]
The identities $q d_E=d_Fq$ and $q(ab)=q(a)q(b)$ prove that $B(q)$ commutes with $b_1$ and $b_2$ separately.

Filter the cone of $B(p)$ by tensor length.
Its length-$n$ quotient is the cone of $(\bar p[1])^{\otimes n}$ with the internal differentials.
The splitting $E=\mathbb C\oplus\bar E$ shows that $\bar p$ is a quasi-isomorphism.
Tensor products of complexes of vector spaces preserve quasi-isomorphisms.
One proof splits every acyclic complex into two-term contractible summands and tensors its contraction with the identity.
Consequently each length quotient of the cone is acyclic.
A cone cycle has finite maximal length.
Solve its boundary equation in that highest quotient and subtract the resulting boundary.
This lowers its maximal length.
Induction ends at length zero and proves that the whole cone is acyclic.
The same argument applies to $B(i)$, and functoriality gives $B(p)B(i)=1$.
\end{proof}

The bar argument uses finite tensor words and polynomial primitives.
A product completion, a vertex algebra, and a chiral sheaf each require a separate comparison.
Even on $A$, the projection $p$ does not preserve the original Poisson brackets:
\[
 p\{h,f\}=-2u,\qquad \{p(h),p(f)\}=0.
\]
Hence this chain algebra map cannot transport a Poisson or Batalin--Vilkovisky operation merely by forgetting its source.

\subsection{Reduction of a specified double complex}

Two odd differentials require their anticommutator to vanish on their common domain.
Their separate nilpotence does not imply this equation.
For example, on $\Lambda(x,y)$, exterior multiplication $d_1=x\wedge-$ and left differentiation $d_2=\partial_x$ are odd and square zero, but
$d_1d_2+d_2d_1=1$.

\begin{theorem}[A comparison under a second differential]
\label{thm:bar-semi-infinite-w}
Let $X^{p,q},Y^{p,q}$ be bigraded complex vector spaces, zero unless $0\leq p\leq N$.
Let $d$ have bidegree $(1,0)$ and $\delta$ have bidegree $(0,1)$, with
\[
 d^2=\delta^2=0,\qquad d\delta+\delta d=0.
\]
Suppose $F:X\to Y$ has bidegree $(0,0)$, commutes with both differentials, and induces a quasi-isomorphism on every column $(X^{p,*},\delta)$.
Then
\begin{equation}\label{eq:bar-semi-infinite-w}
 \operatorname{Tot}(F):
 \left(\bigoplus_{p+q=*}X^{p,q},d+\delta\right)
 \longrightarrow
 \left(\bigoplus_{p+q=*}Y^{p,q},d+\delta\right)
\end{equation}
is a quasi-isomorphism.
The formula is $\operatorname{Tot}(F)((x^{p,q}))=(F(x^{p,q}))$.
\end{theorem}

\begin{proof}
Filter the total complexes decreasingly by the first degree $p$.
This filtration is finite, and the differential on each quotient is $\delta$.
The same filtration on the mapping cone has acyclic quotients by hypothesis.
The short exact sequence of two successive filtration stages gives a long exact sequence in cohomology.
Induction through the finite filtration proves that the cone is acyclic.
The two commutation identities give the total chain identity before any cohomology is taken.
\end{proof}

For quantum Drinfeld--Sokolov reduction, the input category in Arakawa's vanishing theorem is $\mathrm{KL}_k$.
It consists of the specified graded affine modules with locally finite action of the finite-dimensional Lie algebra.
The theorem gives $H_f^j(M)=0$ for $j\ne0$ and exactness of $H_f^0$ on this category \cite[\S6, p.~585; Theorem~7.1(i), p.~588]{Ara15}.
Applying it to a collision complex requires an affine action on each term, membership in this category, and equivariance of all collision maps.
To use the displayed totalization, one also needs its boundedness hypothesis, or a proved substitute for convergence.
The nonlinear $\mathcal W$-algebra and a conformal grading alone do not specify a second semi-infinite differential.

The polynomial maps~\eqref{eq:cbc26-ds-inclusion}--\eqref{eq:cbc26-ds-projection} construct the associative comparison for the classical $\mathfrak{sl}_2$ constraint.
A chiral comparison additionally needs maps from its collision sheaves, including the vacuum terms of every operator product.
The augmentation obstruction in Proposition~\ref{prop:vir-no-augmentation} excludes the augmentation-kernel construction for the matter--ghost vertex algebra.

\subsection{The Virasoro target at a reduced level}

\begin{corollary}[Nilpotence on the Virasoro vacuum carrier]
\label{cor:virasoro-semi-infinite}
Let $k\in\mathbb C\setminus\{-2\}$ and let $M_{c(k)}$ be the universal Virasoro vacuum module with
\[
 c(k)=1-\frac{6(k+1)^2}{k+2}.
\]
Use the ghost module and charge $Q$ of Theorem~\ref{thm:vir-square} on $M_{c(k)}\otimes\mathcal F$.
This charge squares to zero on the full algebraic carrier if and only if
\[
 k=-\frac83\quad\hbox{or}\quad k=-\frac72.
\]
At either level, its absolute vacuum cohomology has one-dimensional components in degrees $0,3$.
Its relative vacuum cohomology has one-dimensional components in degrees $0,2$.
All other components vanish.
\end{corollary}

\begin{proof}
Theorem~\ref{thm:vir-square} gives
\[
 Q^2=\frac{c(k)-26}{12}\sum_{n>0}(n^3-n)c_{-n}c_n.
\]
On $\mathbf1\otimes b_{-2}\Omega$ the output is
$(c(k)-26)\mathbf1\otimes c_{-2}\Omega/2$, a nonzero vector when $c(k)\ne26$.
The equation $c(k)=26$ is $6k^2+37k+56=(3k+8)(2k+7)=0$.
At these levels the matter module is the universal vacuum module of central charge $26$.
The four-state calculation and the nonzero-weight contraction in Proposition~\ref{prop:vir-relative-target-obstruction} compute its absolute and relative cohomology.
\end{proof}

Thus the full displayed ghost carrier cannot define semi-infinite cohomology at arbitrary $k\ne-2$.
Another semi-infinite construction would need another explicitly defined differential and a comparison with this one at the applicable parameters.
The vacuum PBW character $\prod_{n\geq2}(1-q^n)^{-1}$ counts states by conformal weight.
It does not compute the cohomology in the preceding corollary.

\section{Central coefficients and maps over curves}

\subsection{Level reflection and central charge}

\begin{proposition}[Affine scalar normalizations]
\label{cor:anomaly-duality-km}
Fix nonzero $h$ and a positive integer $d$.
For $k\ne-h$, set
\[
 k'=-k-2h,\qquad
 a(k)=\frac{(k+h)d}{2h},\qquad
 c_{\mathrm{Sug}}(k)=\frac{kd}{k+h}.
\]
Then
\[
 a(k)+a(k')=0,\qquad
 c_{\mathrm{Sug}}(k)+c_{\mathrm{Sug}}(k')=2d.
\]
For an affine Lie algebra one takes $h=h^\vee$ and $d=\dim\mathfrak g$.
The half-central-charge normalization has sum $d$, rather than zero.
\end{proposition}

\begin{proof}
Write $t=k+h\ne0$.
Reflection sends $t$ to $-t$.
The two expressions are $a=td/(2h)$ and $c_{\mathrm{Sug}}=d-hd/t$.
Their reflected sums give the formulas.
\end{proof}

\begin{proposition}[Principal central-charge reflection]
\label{cor:anomaly-duality-w}
Fix rank $r$ and vectors $\rho,\rho^\vee$ in a real inner-product space.
For $t\in\mathbb C^*$ define
\[
 c(t)=r-12\left(t\lVert\rho^\vee\rVert^2
             -2(\rho,\rho^\vee)+t^{-1}\lVert\rho\rVert^2\right).
\]
This is the principal central-charge normalization with $t=k+h^\vee$ and long roots of square length $2$.
For a fixed scalar $v$, set $\kappa(t)=v c(t)$.
Then
\[
 c(t)+c(-t)=2r+48(\rho,\rho^\vee),\qquad
 \kappa(t)+\kappa(-t)=v\bigl(2r+48(\rho,\rho^\vee)\bigr).
\]
For $\mathfrak{sl}_2$, $c(t)=13-6(t+t^{-1})$ and $v=1/2$, so the sums are $26$ and $13$.
For $\mathfrak{sl}_3$, $c(t)=50-24(t+t^{-1})$ and $v=5/6$, so the sums are $100$ and $250/3$.
\end{proposition}

\begin{proof}
The $t$ and $t^{-1}$ terms cancel under reflection.
For $\mathfrak{sl}_2$, the positive root $\alpha$ satisfies $(\alpha,\alpha)=2$, and $\rho=\rho^\vee=\alpha/2$ has square $1/2$.
For $\mathfrak{sl}_3$, with simple roots $\alpha_1,\alpha_2$ of square $2$ and $(\alpha_1,\alpha_2)=-1$, one has $\rho=\rho^\vee=\alpha_1+\alpha_2$ of square $2$.
Substitution proves the two cases.
\end{proof}

The parameter reflection $t\mapsto-t$ differs from Feigin--Frenkel duality, whose dual parameters satisfy $t\,t^{\vee}=1/r^\vee$, where $r^\vee$ is the lacing number \cite[p.~566, footnote~1]{Ara15}.
For example, $k=-8/3$ reflects to $k'=-4/3$ in $\mathfrak{sl}_2$.
Their central charges are $26,0$, and their half-central charges sum to $13$.
String ghosts contribute central charge $-26$, or half-central charge $-13$.
They cancel the reflected pair's total central charge.
For a single matter module, Theorem~\ref{thm:vir-square} instead gives the operator coefficient $(C-26)/12$.
These are distinct scalar normalizations attached to specified expressions.

\subsection{The equation for transport of an anomaly operator}

\begin{proposition}[Transport of a square]
\label{prop:cbc26-square-transport}
Let $d_X,d_Y$ be degree-one operators on graded vector spaces and $F:X\to Y$ a degree-zero map.
If $d_YF=Fd_X$, then
\[
 d_Y^2F=Fd_X^2.
\]
In particular, a chain isomorphism between square-zero complexes supplies no identification with a nonzero anomaly operator on another carrier.
If $F$ is injective and $d_Y^2=0$, its intertwining equation forces $d_X^2=0$.
\end{proposition}

\begin{proof}
Apply $d_Y$ to the intertwining equation and use that equation once more.
For an injective $F$, the resulting identity $Fd_X^2=0$ implies $d_X^2=0$.
\end{proof}

\begin{proposition}[Horizontal comparison]
\label{prop:cbc26-horizontal-comparison}
Let $X$ be a smooth manifold.
Let $(E^\bullet,d_E,\nabla_E)$ and $(F^\bullet,d_F,\nabla_F)$ be bounded complexes of finite-rank complex vector bundles with flat connections.
Assume $d_E,d_F$ commute with the connections.
On $\Omega^p(X,E^q)$ use total differential $\nabla_E+(-1)^p d_E$, and use the same convention for $F$.
A degree-zero bundle map $\phi:E^\bullet\to F^\bullet$ induces a map of these total complexes if and only if
\[
 d_F\phi=\phi d_E,\qquad
 \nabla_F\phi=(1\otimes\phi)\nabla_E.
\]
If these identities hold and $\phi$ is a quasi-isomorphism on each fiber, it induces an isomorphism on the hypercohomology of the two local-system complexes.
\end{proposition}

\begin{proof}
On a section of form degree zero, the two components of the total chain equation lie in different form degrees.
They give exactly the displayed identities.
The Leibniz rule extends these identities to arbitrary forms.
Flat local frames identify each stalk of the kernel local systems with its fiber.
The cone of the induced local-system map is therefore acyclic on every stalk.
Its cohomology sheaves vanish, so it is zero in the derived category of sheaves.
Applying derived global sections gives the asserted hypercohomology isomorphism.
\end{proof}

\begin{example}[A genus-one obstruction to fiberwise transport]
\label{ex:cbc26-torus-monodromy}
On the oriented torus $T^2=\mathbb R^2/\mathbb Z^2$, let $L_\lambda$ be the rank-one complex local system with monodromy $\lambda\in\mathbb C^*$ on the first circle and $1$ on the second.
Its cellular cochain complex, with the product orientation, is
\[
 \left(\Lambda^\bullet\mathbb C^2,
       (\lambda-1)e_1\wedge-\right).
\]
This is the tensor product of the two circle complexes with differentials $\lambda-1$ and $0$.
If $\lambda\ne1$, contraction by the dual vector $e_1^*$, divided by $\lambda-1$, gives $dh+hd=1$.
Thus $H^*(T^2,L_\lambda)=0$.
For $\lambda=1$, the differential is zero and the dimensions are $(1,2,1)$.
The fibers in both cases are one-dimensional complexes concentrated in degree zero.
No fiberwise dimension or scalar central coefficient determines the genus-one cohomology.
\end{example}

A comparison over the moduli of curves needs maps of its sheaves or local systems, compatible with the connections and with clutching at the boundary.
The genus-zero filtered lifting criterion alone does not define such maps.
The same requirement applies separately to pairings, anomaly operators, and the products used in a Maurer--Cartan equation.
